% Generated draft: requires mathematical and source review.
\documentclass[11pt]{article}
\usepackage[margin=25mm]{geometry}
\usepackage{fontspec}
\setmainfont{Latin Modern Roman}
\usepackage{amsmath,amssymb,mathtools,unicode-math}
\setmathfont{Latin Modern Math}
\usepackage{enumitem,needspace,xurl,hyperref}
\setlength{\parindent}{0pt}
\setlength{\parskip}{5pt}
\setlength{\emergencystretch}{3em}
\newcommand{\N}{\mathbb{N}}
\newcommand{\Z}{\mathbb{Z}}
\newcommand{\Q}{\mathbb{Q}}
\newcommand{\R}{\mathbb{R}}
\newcommand{\C}{\mathbb{C}}
\newcommand{\F}{\mathbb{F}}
\newcommand{\A}{\mathbb{A}}
\newcommand{\PP}{\mathbb{P}}
\newcommand{\E}{\mathbb{E}}
\newcommand{\eps}{\varepsilon}
\newcommand{\dd}{\,\mathrm{d}}
\begin{document}
\section*{UnsolvedMath ID: 30001774}
\textbf{Source problem code:} catalogue entry
\subsection{Definitions and mathematical statement}
The supplied record groups three families of questions. Their formulations below are conditional on unspecified representation-theoretic conventions; the record does not provide enough hypotheses to turn them into fully quantified conjectures.

Let \(k\) denote the coefficient field, whose characteristic is not specified, and let \(S_m\) denote the symmetric group on \(m\) letters. Interpret symmetric-group modules as modules over the group algebra \(kS_m\). The permitted values of the integer parameters and any finite-dimensionality assumptions remain unspecified.

\begin{enumerate}
\item \textbf{Young dimension and injective dimension.} For an \(S_n\)-module \(M\) in the intended module category, the question is whether
\[
\operatorname{Ydim}(M)=\operatorname{injdim}_{A}G(M),
\]
where \(A\) is the intended Schur algebra and \(G(M)\) is the associated \(A\)-module. For a specified algebra \(A\), the injective dimension of an \(A\)-module \(N\) is the least integer \(r\geq 0\) for which an injective resolution
\[
0\longrightarrow N\longrightarrow I^0\longrightarrow \cdots\longrightarrow I^r\longrightarrow 0
\]
exists, and is \(\infty\) if no finite such resolution exists, subject to the chosen convention for the zero module. The notation \(\operatorname{Ydim}\) is a placeholder for the source's undefined term ``Young dimension.'' The record does not specify its definition, the parameters of \(A\), the construction or direction of \(G\), or the range of \(M\) for which the equality is asked. These choices cannot be inferred from the displayed equality.

\item \textbf{Self-extension algebras of Young modules.} Let \(d\) be in the intended range, let \(B\) be a block of \(kS_d\), and let \(\lambda\vdash d\) label a Young module \(Y^\lambda\) in \(B\). Here a block means an algebra summand determined by a primitive central idempotent, and \(\lambda\vdash d\) means that \(\lambda\) is a partition of \(d\). The source uses the conventional Young-module notation \(Y^\lambda\), without specifying its construction. For every such \(\lambda\) that is maximal in its block in the intended, unspecified order, form the Yoneda algebra
\[
E_\lambda=\bigoplus_{i\geq 0}\operatorname{Ext}^{i}_{kS_d}(Y^\lambda,Y^\lambda),
\]
whose multiplication is composition of extensions. Is \(E_\lambda\) graded-commutative, meaning that
\[
xy=(-1)^{ij}yx
\qquad
\text{for all }x\in E_\lambda^i,\ y\in E_\lambda^j,\ i,j\geq 0?
\]
The accompanying question asks whether this algebra can act naturally on the modules in \(B\). The record does not define the order used for maximality or what constitutes this natural action. In particular, it does not specify whether the action is on modules themselves, on extension groups, or on objects of a derived category. Consequently, the action question cannot be expressed as a definite existence assertion from the supplied information.

\item \textbf{Specht and dual Specht filtrations.} Write \(S^\mu\) for the Specht module indexed by a partition \(\mu\vdash m\), and write \((S^\mu)^*=\operatorname{Hom}_k(S^\mu,k)\) for its contragredient dual. A Specht filtration of a \(kS_m\)-module \(V\) is a finite chain of submodules
\[
0=V_0\subseteq V_1\subseteq\cdots\subseteq V_r=V
\]
whose successive quotients are Specht modules. A dual Specht filtration has successive quotients of the form \((S^\mu)^*\).

Using the conventional interpretation, \(\operatorname{Lie}(n)\) is the multilinear component of the free Lie algebra over \(k\) on \(n\) generators, with \(S_n\) acting by permutation of the generators. The record asks whether \(\operatorname{Lie}(n)\) admits a Specht or dual Specht filtration, and asks the analogous question for ``the Foulkes module.'' Neither the field restrictions nor the integer ranges are supplied. The Foulkes module's construction, parameters, ambient symmetric group, and any twisting convention are also unspecified. Finally, the wording does not settle whether the source asks for at least one of the two filtration types or asks separately about each type. These alternatives remain unresolved.
\end{enumerate}
\subsection{Short English statement}
The record asks about equality of Young and injective dimensions, graded commutativity and natural actions of certain Young-module extension algebras, and Specht or dual Specht filtrations of Lie and Foulkes modules. Missing field hypotheses, functor definitions, ordering conventions, and module parameters prevent a fully precise joint formulation. These are statement-only questions requiring mathematical review.
\subsection{Sources}
\begin{itemize}
\item Local UnsolvedMath registry, supplied record with permanent ID \(30001774\), source problem code \texttt{catalogue entry}, titled ``Homological and Filtration Problems for Symmetric-Group Modules,'' in the collection ``Oberwolfach Reports Open Problems.'' No individual proposer is identified in the supplied record; the proposed year is 2011.
\item Supplied report citation: \emph{Mini-Workshop: Modular Representations of Symmetric Groups and Related Objects} (2011), \emph{Oberwolfach Reports}, pp.\ 1177--1205, DOI \texttt{10.4171/owr/2011/21}. Supplied source URL: \url{https://doi.org/10.4171/owr/2011/21}. The background attributes the questions to workshop open-problem discussions documented by this report; no question-specific page or speaker attribution is supplied.
\item Supplied literature source, labelled ``original OWR report'': \url{https://doi.org/10.4171/owr/2011/21}. The supplied literature assessment was checked on \texttt{2026-08-21} and labels the record open, describing it as several unrelated questions about block algebras and Lie and Foulkes modules rather than one resolved claim. This is supplied dated curation, not an independently verified literature assessment. The background explicitly requests independent verification of time-sensitive status and cited claims.
\item The registry labels the statement \texttt{corrected\_verified} and reports a manual semantic, typographic, and LaTeX review. These are supplied editorial claims; they do not resolve the missing mathematical specifications identified here. The cited report has not been independently consulted for this draft.
\end{itemize}
\end{document}
